% !TEX root = ../../../proposal.tex

\subsection{Compositional Diffusion Models}
\label{sec:compositional_diffusion}
Single text-to-image diffusion models often fail when a prompt requires several factors to
be combined correctly. Recent analyses \citep{t2i_compbench_2025, du2024compositional,
mechanisms2025projective} show consistent errors in attribute binding, spatial relations,
and numeracy, even when the image remains visually plausible. These failures arise because
the models largely learn correlations in joint image-text data rather than an explicit
factorized structure over objects, attributes, and relations \citep{t2i_compbench_2025}.
In the proposed study, we refer to the conventional case in which a single natural-language prompt
specifies all desired factors and the model attempts to realize them jointly as
\emph{semantic composition}.

The next three definitions review the main formulations of composition in this literature: the \emph{simple product}, \emph{Bayes composition}, and \emph{projective composition}. 
The first two define composition by combining densities directly, while the third defines it through what the composed result should preserve.

\subsubsection{Simple Product}

The first definition is the \emph{simple product}. Given two distributions \(p_1\) and
\(p_2\) over \(\mathbb{R}^n\), it defines the composition as
\[
    \hat p(x) \propto p_1(x)p_2(x).
\]
This is a familiar way to combine distributions, but it has an important limitation: the
composed distribution cannot be truly out-of-distribution with respect to \(p_1\) or
\(p_2\). If either \(p_1(x)=0\) or \(p_2(x)=0\), then \(\hat p(x)=0\) as well. As
\citet{mechanisms2025projective} show, this becomes a problem when the desired
composition lies outside the support of both source distributions. In their CLEVR example,
if each source distribution contains only one object, then any image with \emph{both}
objects has zero probability under both \(p_1\) and \(p_2\), so the simple product also
assigns it zero probability. 

\subsubsection{Bayes Composition}

A second prior definition is \emph{Bayes composition}. 
Rather than multiplying arbitrary densities, Bayes composition treats the component
distributions as conditionals of a common underlying data distribution.~\citep{liu2022compositional,du2023reduce,du2024compositional,NEURIPS2020_49856ed4}.
If \(c_1\) and \(c_2\) denote two conditions of interest, the
desired composition is defined as
\[
    \hat p(x) := p(x \mid c_1, c_2).
\]
Assuming \(c_1\) and \(c_2\) are conditionally independent given \(x\), Bayes' rule gives
\begin{equation}
    \hat p(x)
    = p(x \mid c_1, c_2)
    \propto
    \frac{p(x \mid c_1)\,p(x \mid c_2)}{p(x)}.
    \label{eq:poe_dist}
\end{equation}
This formulation does not define a truly out-of-distribution composition
\citep{mechanisms2025projective}, because it works when the desired conjunction already corresponds
to a meaningful conditional under the training distribution, but cannot in general create a
new conjunction that lies outside that support.

For diffusion models, this composition is implemented at the level of the score.
Taking \(\nabla_x \log\) of Equation~\ref{eq:poe_dist} gives
\begin{equation}
    \nabla_x \log \hat p(x \mid c_1, c_2)
    =
    \nabla_x \log p(x \mid c_1)
    +
    \nabla_x \log p(x \mid c_2)
    -
    \nabla_x \log p(x).
    \label{eq:poe_score}
\end{equation}
Thus, the composed score is the sum of the two conditional scores minus the unconditional
score.

Using the score estimate in Equation~\ref{eq:score_noise}, the score
composition in Equation~\ref{eq:poe_score} can be written directly in terms of noise predictions. Let
\(\varepsilon_\theta(x_t,t,\varnothing)\) denote the unconditional prediction and
\(\varepsilon_\theta(x_t,t,c_i)\) the prediction conditioned on concept \(c_i\). For a
single condition, classifier-free guidance takes the form
\[
    \varepsilon_\theta(x_t,t,\varnothing)
    +
    w_i\bigl[\varepsilon_\theta(x_t,t,c_i)-\varepsilon_\theta(x_t,t,\varnothing)\bigr].
\]
Applying this construction to both concepts yields the practical AND operator used by
\citet{liu2022compositional}:
\begin{equation}
    \tilde\varepsilon_{\mathrm{AND}}
    =
    \varepsilon_\theta(x_t,t,\varnothing)
    +
    w_1\bigl[\varepsilon_\theta(x_t,t,c_1)-\varepsilon_\theta(x_t,t,\varnothing)\bigr]
    +
    w_2\bigl[\varepsilon_\theta(x_t,t,c_2)-\varepsilon_\theta(x_t,t,\varnothing)\bigr].
    \label{eq:and_operator}
\end{equation}
When \(w_1=w_2=1\), this reduces to
\[
    \varepsilon_\theta(x_t,t,c_1)
    +
    \varepsilon_\theta(x_t,t,c_2)
    -
    \varepsilon_\theta(x_t,t,\varnothing),
\]
which is exactly the Bayes composition written in noise-prediction form.

\subsubsection{Projective Composition}
\label{subsec:bg_projective}

\citet{mechanisms2025projective} argue that a useful definition of composition must answer
two separate questions: what the composed distribution should preserve, and when a practical
compositional method can realize it. Their notion of \emph{projective composition} addresses
the first question. Under this view, the composed distribution is judged by the part each
source is meant to contribute. For example, if one source determines style and another
content, then a composition is correct if it preserves the style of the first and the
content of the second, even when neither source distribution contains that full combination.
This is closely aligned with \citet{du2024compositional}, where composition is understood as
building a new distribution from simpler components that each contribute part of the final
structure.

This matters because it allows composition that is out of distribution with respect to the
source distributions being combined. The final sample can combine properties from different
sources without needing to belong to any one source distribution as a whole. 
As Figure~\ref{fig:projective_composition} illustrates, the composed distribution can differ
from both source distributions in the original sample space while still matching each source
after applying the corresponding projection.
That is precisely what the simple product cannot express, since it is confined to overlapping
support, and what Bayes composition cannot justify when the desired conjunction is not
already supported by some underlying joint distribution.

\begin{figure}[H]
    \centering
    \includegraphics[width=0.45\linewidth]{chapters/background/media/projective_composition.png}
    \caption{Illustration of projective composition. The composed distribution \(\hat p\) need not coincide with either source distribution \(p_1\) or \(p_2\) in the original sample space. Instead, correctness is defined after applying source-specific projections: under \(\Pi_1\), the composed distribution matches the aspect preserved from \(p_1\), and under \(\Pi_2\), it matches the aspect preserved from \(p_2\). The figure therefore shows how a composition can be out of distribution with respect to the sources as whole distributions while still preserving the designated contribution of each source.}
    \label{fig:projective_composition}
\end{figure}

The key implication is representational. 
Standard disentanglement is not enough if concepts still interfere when one is changed. What composition requires is a stronger separation in which each concept acts on its own factor without perturbing the others. 
\citet{mechanisms2025projective} further show that this separation need not appear in pixel space itself. 
If it exists in some invertible latent space, then the same composition operator yields a correct projective composition with respect to those latent space factors. 
This is a statement about when the target composition is well defined, not a guarantee that reverse-diffusion sampling will always recover it. 
For the proposed study, the lesson is that compositional success depends on whether the learned representation allows the concepts to be preserved independently.
